\section{Compute nobody's customer controls}
\label{sec:27.3}
The second piece is compute bought by a body that buys and judges nothing. It procures time on machines and allocates it as the panels of \S\ref{sec:27.5} decide; it grants no finding, holds no placement and funds no particular answer, and it publishes every allocation.

The model is CERN. Its member states founded it by a convention in 1953 \autocite{cern1953convention}, and the twenty-three of them fund it in proportion to their net national income \autocite{cern2026participation}, so the laboratory belongs to no one of them.

The American start is the National Artificial Intelligence Research Resource. The National Science Foundation launched it as a pilot in January 2024, with fourteen other federal agencies and twenty-eight private-sector partners; in its first eighteen months it connected more than four hundred research teams and educators to computing and data, and in 2025 the Foundation sought an operations center to carry it toward a long-term resource \autocite{nsf2025nairroc}. Two things about it are worth keeping in view. It is a pilot, which a later budget can end without repealing anything. And part of what it lends is lent by private partners, who can stop lending.

The obvious American host for public compute is the Department of Energy, and that isn't simple. Its Office of Science runs open user facilities, Frontier at Oak Ridge among them, and gives most of their time to outside researchers through a competitive call \autocite{ornl2023incite}. The same department's National Nuclear Security Administration bought El Capitan, at Lawrence Livermore, to maintain the nuclear weapons stockpile \autocite{nnsa2019elcapitan}. So the separation Appendix~\ref{sec:A} draws for the payer is drawn for the buyer too: the procurement body sits in no department that operates a deployment, may buy time on open facilities, and buys none through a weapons program's chain of allocation.

Beside it sit the arrangements the book already has: the cooperative facility of \S\ref{sec:16.2}, which a county or a cooperative owns, and the levy of chapter~\ref{sec:33}, which frees the computing, energy and hardware that public uses would otherwise bid away.
